\documentclass[article,nojss]{jss}

%% --- arXiv-first preprint. The `nojss` option drops the journal masthead and
%% logo, so only jss.cls and jss.bst are needed (no jsslogo.jpg). Obtain them from
%% https://www.jstatsoft.org/style; arXiv's TeXLive also carries the jss package.
%% For an eventual JSS journal submission, remove the `nojss` option.

\usepackage{amsmath,amssymb,booktabs}

%% --- author / title metadata ----------------------------------------------

\author{Neal Caren\\University of North Carolina at Chapel Hill}
\Plainauthor{Neal Caren}

\title{\pkg{topica}: A Unified \proglang{Python} Framework for Topic Modeling
in the Social Sciences}
\Plaintitle{topica: A Unified Python Framework for Topic Modeling in the Social
Sciences}
\Shorttitle{\pkg{topica}: A Unified Framework for Topic Modeling}

\Abstract{
Topic modeling is now a standard tool for social scientists who work with text,
but the models they need are scattered across incompatible software: the
structural topic model lives in \proglang{R}, the fastest collapsed-Gibbs
sampler lives in \proglang{Java}, and the keyword-assisted and embedding-based
models live in separate \proglang{Python} repositories, none sharing a data
format or a diagnostic suite. We present \pkg{topica}, a \proglang{Python}
framework that brings more than a dozen topic models---classical and
embedding-based alike---behind one \pkg{NumPy}-native interface. Every model
exposes the same surface, the topic-word and document-topic matrices, so one set
of coherence, exclusivity, stability, labeling, and covariate-effect tools
applies to all of them, and a count-based model and a clustering of embeddings
can be compared on a single corpus in one session. That same uniformity is what
makes the framework extensible: a new model that presents the same surface
inherits the entire diagnostic and effect-estimation stack without writing any
of it. \pkg{topica} is built around the premise that the researcher, not the
software, owns the theoretical decisions that make a topic-model study
credible---what the corpus represents, how many topics is the right granularity,
what each topic means---and supplies the mechanics and the honest diagnostics to
support them, declining, for instance, to report confidence intervals for a
model that has no posterior to draw them from. A parallel \proglang{Rust} core
makes the fits fast---three to twenty-two times faster than \proglang{R} \pkg{stm}
for the structural and other variational models, at parity with the compiled
\pkg{MALLET} and \pkg{keyATM} samplers---and reproducible, the variational models
to the bit and the samplers from a fixed seed, and each model is validated against
its reference implementation. We describe the design, the
model family, and a principled analysis workflow.
}

\Keywords{topic models, structural topic model, computational social science,
text as data, \proglang{Rust}, \proglang{Python}, reproducibility,
\pkg{topica}}
\Plainkeywords{topic models, structural topic model, computational social
science, text as data, Rust, Python, reproducibility, topica}

\Address{
  Neal Caren\\
  Department of Sociology\\
  University of North Carolina at Chapel Hill\\
  E-mail: \email{neal.caren@unc.edu}\\
  URL: \url{https://nealcaren.github.io/topica/}
}

\begin{document}

%% ===========================================================================
\section{Introduction}\label{sec:intro}
%% ===========================================================================

Text is now ordinary data in the social sciences. Sociologists read topic
structure off newspaper archives, political scientists measure how the framing
of an issue shifts across parties, and survey researchers code open-ended
responses at a scale that hand-coding cannot reach \citep{dimaggio2013exploiting,
grimmer2013text, roberts2014structural, gentzkow2019text, grimmer2022text}. The
workhorse for this kind of analysis is the probabilistic topic model
\citep{blei2003lda, blei2012probabilistic}, and over two decades the family has
grown to cover correlated topics, topics that change over time, topics that
depend on document metadata, topics steered by a researcher's keywords, and
topics read off neural embeddings.

The software has not kept pace with the breadth of the methods. A social
scientist who wants the structural topic model, the field standard for relating
topic prevalence to covariates, must work in \proglang{R} with the \pkg{stm}
package \citep{roberts2016model, roberts2019stm}. The fastest collapsed-Gibbs
sampler ships in \proglang{Java} as \pkg{MALLET} \citep{mccallum2002mallet}. The
keyword-assisted model \citep{eshima2024keyatm} is a separate \proglang{R}
package, and the embedding-based models \citep{grootendorst2022bertopic,
angelov2020top2vec} are separate \proglang{Python} repositories again. Each tool
defines its own document format, returns topics in its own structure, and
computes its own subset of diagnostics, if any. A study that compares two model
families therefore spans two languages and three data conventions, leaving the
analyst to glue the pieces together by hand.

Three problems follow from this fragmentation. First, comparison is expensive:
asking whether the structural topic model and a clustering of embeddings find
the same themes in a corpus means running two pipelines that do not speak to each
other and aligning the output by hand. Second, diagnostics are uneven: coherence
and exclusivity are built into some tools and absent from others, so the standard
of evidence depends on which package produced the result. Third, and
most consequential for cumulative social science, reproducibility is not
guaranteed. Several reference implementations do not fix all sources of
randomness, so the same script on the same data can return different topics on a
second run, which undermines the replication that the field increasingly
expects.

We present \pkg{topica}, a \proglang{Python} library that addresses these three
problems together. \pkg{topica} provides more than a dozen topic models behind a
single \pkg{NumPy}-native interface, computes the standard diagnostics for any of
them, and produces reproducible fits. It rests on four claims, which the rest of
this article defends in turn:

\begin{enumerate}
\item \textbf{Breadth in one interface.} Count-based and embedding-based models
share one contract, so the same diagnostic, labeling, and covariate-effect tools
apply to every model and cross-model comparison costs no extra plumbing
(Sections~\ref{sec:design} and~\ref{sec:models}).
\item \textbf{Speed.} A parallel \proglang{Rust} core fits the structural and
other variational models from three to twenty-two times faster than \proglang{R}
\pkg{stm}, and matches the compiled \pkg{MALLET} and \pkg{keyATM} samplers core
for core---without a JVM and without \pkg{PyTorch} (Section~\ref{sec:performance}).
\item \textbf{Reproducibility by construction.} The variational models are
deterministic to the bit even when multithreaded, and the sampling models are
seed-reproducible (Section~\ref{sec:design}).
\item \textbf{Reference validation.} Each model is checked against its original
implementation, and we report the agreement honestly, including where two correct
engines disagree because the model is non-convex (Section~\ref{sec:validation}).
\end{enumerate}

A fifth commitment runs through the design rather than appearing as a feature.
A topic model is a measurement instrument, not a discovery oracle, and the
decisions that make a topic-model study credible are theoretical: what the corpus
represents, how many topics is the right granularity for the question, what each
topic means, which covariates belong in the model, and whether a pattern is
substantively interesting. \pkg{topica} is built so that the researcher makes
those decisions and the software supplies the mechanics and the honest
diagnostics to support them (Section~\ref{sec:workflow}). The library refuses to
manufacture certainty the data does not support: it declines to report confidence
intervals for a model that has no posterior to draw them from, rather than
returning a plausible-looking number.

%% ===========================================================================
\section{Related software}\label{sec:related}
%% ===========================================================================

Topic modeling is well served by individual tools, each strong in its niche.
\pkg{MALLET} \citep{mccallum2002mallet} provides a fast, well-tested
collapsed-Gibbs sampler for LDA and its Dirichlet-multinomial regression
extension \citep{mimno2008dmr}, with hyperparameter optimization, in
\proglang{Java}. \pkg{gensim} \citep{rehurek2010gensim} offers online LDA and a
dynamic topic model in pure \proglang{Python}. \pkg{tomotopy} \citep{tomotopy}
implements many models with a fast \proglang{C++} core. In \proglang{R}, the
\pkg{stm} package \citep{roberts2019stm} is the reference for the structural
topic model, with the diagnostic, effect-estimation, and visualization tooling
that social scientists have standardized on, and \pkg{keyATM}
\citep{eshima2024keyatm} provides keyword-assisted models. On the neural side,
\pkg{BERTopic} \citep{grootendorst2022bertopic} and \pkg{Top2Vec}
\citep{angelov2020top2vec} cluster document embeddings, and \pkg{ETM}
\citep{dieng2020etm} and \pkg{FASTopic} \citep{wu2024fastopic} are generative
embedding models.

What no single tool offers is the combination a social-science workflow needs:
the structural topic model and the fast classical samplers and the embedding
models, behind one interface, with shared diagnostics, reproducible fits, and
honest uncertainty for covariate effects. Table~\ref{tab:compare} sketches the
gap. The point is not that any one tool is deficient in its own domain; it is
that no row but the last spans the matrix, so an analysis that needs more than
one cell pays a translation cost at every boundary.

\begin{table}[t!]
\centering
\caption{Topic-modeling software and the social-science feature matrix, as of the
current released versions. ``STM/content'' is a full structural topic model with
content covariates; \emph{partial} marks a prevalence-only covariate model
(Dirichlet-multinomial regression or its keyword analogue), which does not let the
topic-word distribution vary by covariate. ``Effect SEs'' marks built-in
covariate-effect estimation with uncertainty. ``Determinism'' marks the documented
reproducibility guarantee: \emph{seed} reproduces from a fixed seed,
\emph{partial} has a non-seedable stage (UMAP in \pkg{BERTopic}), and a check marks
\pkg{topica}'s bit-for-bit variational fits.}
\label{tab:compare}
\small
\begin{tabular}{lccccccc}
\toprule
 & Language & Classical & STM/content & Short-text & Embedding & Effect SEs & Determinism \\
\midrule
\pkg{MALLET}    & \proglang{Java}    & \checkmark & partial    & --         & --         & --         & seed \\
\pkg{gensim}    & \proglang{Python}  & \checkmark & --         & --         & --         & --         & seed \\
\pkg{tomotopy}  & \proglang{Python}/\proglang{C++} & \checkmark & partial & \checkmark & --      & --         & seed \\
\pkg{stm}       & \proglang{R}       & \checkmark & \checkmark & --         & --         & \checkmark & seed \\
\pkg{keyATM}    & \proglang{R}       & \checkmark & partial    & --         & --         & \checkmark & seed \\
\pkg{BERTopic}  & \proglang{Python}  & --         & --         & \checkmark & \checkmark & --         & partial \\
\midrule
\pkg{topica}    & \proglang{Python}/\proglang{Rust} & \checkmark & \checkmark & \checkmark & \checkmark & \checkmark & \checkmark \\
\bottomrule
\end{tabular}
\end{table}

%% ===========================================================================
\section{Design and architecture}\label{sec:design}
%% ===========================================================================

\pkg{topica} is two layers: a \proglang{Rust} core that does the numerical work
and a thin \proglang{Python} layer that presents it. The core is compiled to a
native extension module with \pkg{PyO3} \citep{pyo3} and \pkg{maturin}
\citep{maturin}, and built against the stable ABI so one wheel runs on every
supported \proglang{Python} version. The choice of \proglang{Rust} follows from
the workload: collapsed-Gibbs sampling and variational EM are tight loops over
millions of tokens, which a compiled language runs an order of magnitude faster
than interpreted \proglang{Python}, and \proglang{Rust}'s thread model lets the
core parallelize across cores without the global interpreter lock.

\subsection{One contract for every model}

The boundary between the two layers is deliberately narrow: every model takes a
corpus to \code{fit}, and afterward exposes the two matrices that define a topic
model, the topic-word distribution $\phi$ as \code{topic\_word} and the
document-topic distribution $\theta$ as \code{doc\_topic}, both as plain
\pkg{NumPy} arrays \citep{harris2020numpy}, together with \code{top\_words} and
\code{save}/\code{load}. Listing the interface is short enough to give in full:

\begin{CodeChunk}
\begin{CodeInput}
import topica

model = topica.LDA(num_topics=20, seed=42)
model.fit(docs, iters=1000)

model.topic_word        # phi: ndarray (num_topics, vocab_size)
model.doc_topic         # theta: ndarray (num_docs, num_topics)
model.top_words(10)     # the ten highest-probability words per topic
\end{CodeInput}
\end{CodeChunk}

Because the diagnostics consume only \code{topic\_word} and \code{doc\_topic},
we write them once and they apply to every model. \code{topica.coherence},
\code{topica.exclusivity}, \code{topica.estimate\_effect}, and the rest do not
know or care which model produced the matrices they are given. This is what makes
cross-model comparison free: an LDA fit and a clustering of embeddings present
the same surface, so the analyst can score both with the same coherence measure
and align their topics with the same tools.

\subsection{Extensibility falls out of the contract}

Because the analysis surface reads only from \code{topic\_word},
\code{doc\_topic}, \code{vocabulary}, and \code{top\_words}, adding a model is a
local act, and so is bringing one \pkg{topica} does not ship. Any object that
presents those four members---whether a built-in model, a subclass, or an
external estimator the analyst wraps---is scored, labeled, and compared by the
same functions, with no code written for it on the diagnostics side:

\begin{CodeChunk}
\begin{CodeInput}
class MyModel:
    def fit(self, docs): ...   # produce the two matrices from a corpus
    topic_word                 # ndarray (num_topics, vocab_size)
    doc_topic                  # ndarray (num_docs, num_topics), rows sum to 1
    vocabulary                 # list[str], column order of topic_word
    def top_words(self, n=10, *, topic=None): ...

m = MyModel(); m.fit(docs)
topica.coherence(m, docs)          # every diagnostic ...
topica.exclusivity(m)
topica.topic_diversity(m)
topica.estimate_effect(m.doc_topic, X)   # ... and effect tool applies to m.
\end{CodeInput}
\end{CodeChunk}

The built-in models are written to exactly this contract; the implementer's guide
that ships with the package documents it for contributors.

\subsection{Determinism as a constraint}

We engineer reproducibility rather than hope for it. The
variational models (the correlated and structural topic models, the dynamic topic
model, and supervised LDA) produce bit-for-bit identical fits on repeated runs,
including when the variational E-step is parallelized across cores, because the
parallel reductions are performed in a fixed order rather than in whatever order
the threads happen to finish. The collapsed-Gibbs models are reproducible from a
fixed seed. The practical consequence is that the most basic robustness check a
topic-model study should report, refitting under different seeds to confirm that
the topics persist, is a clean experiment in \pkg{topica}: the only thing that
varies is the seed.

\subsection{A lean core with optional extras}

The core depends only on \pkg{NumPy}. Features that pull heavier dependencies are
optional extras, so the base install stays small and \pkg{PyTorch} is never
required, even for the embedding-based models, which take vectors the user
supplies from any embedder, and even for the neural \code{ProdLDA}, whose
amortized variational autoencoder is hand-coded against \pkg{NumPy} alone. The \code{viz} extra adds the plotting stack, the
\code{formula} extra adds \proglang{R}-style model formulas, and the \code{llm}
extra adds optional large-language-model labeling and embedding through an API or
a local model. The embedding models build on two pure-\proglang{Rust},
BLAS-free crates for HDBSCAN \citep{campello2013hdbscan} and the optional UMAP
reducer \citep{mcinnes2018umap}.

%% ===========================================================================
\section{The model family}\label{sec:models}
%% ===========================================================================

\pkg{topica} groups its models into two kinds by how they treat words.
\emph{Count-based} models learn topics from word counts, by collapsed-Gibbs
sampling, variational EM, or an amortized variational autoencoder, and present
topic-word weights on the probability simplex. \emph{Embedding-based} models
start from document vectors the user supplies. Table~\ref{tab:models} lists the family and the question each model
answers. The breadth is the point: a single import covers the span from the LDA
baseline to the structural topic model to neural clustering.

\begin{table}[t!]
\centering
\caption{The \pkg{topica} model family. All share the \code{fit}\,$\rightarrow$\,
\code{topic\_word}/\code{doc\_topic} contract of Section~\ref{sec:design}.}
\label{tab:models}
\small
\begin{tabular}{lll}
\toprule
Model & What it is for & Reference \\
\midrule
\multicolumn{3}{l}{\emph{Count-based}}\\
\code{LDA}          & Classic topics, fast collapsed-Gibbs (SparseLDA)     & \citet{blei2003lda, yao2009sparselda} \\
\code{DMR}          & Topics conditioned on document metadata              & \citet{mimno2008dmr} \\
\code{LabeledLDA}   & Supervised topics tied to document labels            & \citet{ramage2009llda} \\
\code{CTM}          & Correlated topics (logistic-normal)                  & \citet{blei2007ctm} \\
\code{ProdLDA}      & Sharper topics via a product-of-experts VAE          & \citet{srivastava2017prodlda} \\
\code{STM}          & Prevalence \emph{and} content covariates             & \citet{roberts2016model} \\
\code{SAGE}         & A topic worded differently across groups             & \citet{eisenstein2011sage} \\
\code{HDP}          & Infers the number of topics                          & \citet{teh2006hdp} \\
\code{DTM}          & Topics that evolve across time slices                & \citet{blei2006dtm} \\
\code{SupervisedLDA}& Topics shaped to predict a response                  & \citet{blei2008slda} \\
\code{PT}, \code{GSDMM} & Short-text models                                & \citet{zuo2016pt, yin2014gsdmm} \\
\code{SeededLDA}, \code{KeyATM} & Guided topics from seed words            & \citet{jagarlamudi2012seeded, eshima2024keyatm} \\
\code{PA}, \code{HLDA} & Topic hierarchies                                 & \citet{li2006pa, griffiths2004hlda} \\
\midrule
\multicolumn{3}{l}{\emph{Embedding-based}}\\
\code{BERTopic}     & Cluster embeddings, label by class-TF-IDF            & \citet{grootendorst2022bertopic} \\
\code{Top2Vec}      & Topics as points in the embedding space              & \citet{angelov2020top2vec} \\
\code{ETM}          & Generative LDA factored through embeddings           & \citet{dieng2020etm} \\
\code{FASTopic}     & Topics from optimal-transport plans                  & \citet{wu2024fastopic} \\
\bottomrule
\end{tabular}
\end{table}

Two implementations carry most of the weight for social-science users and are
worth singling out. The structural topic model is implemented in full, with
topic prevalence and topic content both allowed to depend on covariates through
the SAGE mechanism, so the same topic can be worded differently across groups
while its prevalence is regressed on metadata. The keyword-assisted model is also
implemented in full, covering the base, covariate, and dynamic variants, the
information-theoretic token weighting, and the change-point hidden Markov model
\citep{chib1998changepoint} of the dynamic specification.

%% ===========================================================================
\section{A principled analysis workflow}\label{sec:workflow}
%% ===========================================================================

The interface in Section~\ref{sec:design} supports a workflow that mirrors how a
careful text-analysis study is built, with the researcher making the
substantive calls and \pkg{topica} supplying the mechanics. The division of labor
between researcher and software is part of the design, so we walk the workflow to
show where each one decides.

\subsection{Build a defensible corpus}

The corpus fixes the scope of every later claim, so preprocessing is a
substantive decision rather than a chore. \pkg{topica} provides the \code{Corpus}
container and tokenization, phrase learning, and document-splitting tools, and the
workflow documents every choice (stopwords, frequency thresholds, phrases) because
the methods section has to report them.

\subsection{Choose and justify the number of topics}

The number of topics $K$ is the decision agents and automated pipelines most
often get wrong, because the diagnostics tempt one to optimize it. $K$ is a
research decision about granularity, not a hyperparameter to maximize: a model
with fewer topics often scores higher on coherence while collapsing distinctions
the researcher cares about. \code{topica.search\_k} fits a range of $K$ and
reports coherence, exclusivity, and held-out perplexity for each, and
\code{topica.quality\_frontier} plots the per-topic coherence-exclusivity
trade-off. The honest procedure is to fit a small range around a theoretically
plausible $K$, label the topics at each, count the uninterpretable ones, and let
the researcher choose, reporting sensitivity to that choice.

\subsection{Fit, read, and validate}

Topics are read from several angles together, since top words alone underdetermine
meaning: \code{label\_topics} reports the highest-probability, FREX
\citep{bischof2012frex}, and lift words, and \code{find\_thoughts} returns the
documents that load most heavily on a topic. Validation is not optional.
\code{topica.coherence} computes the windowed coherence measures
\citep{lau2014coherence} in the core, \code{topica.exclusivity} and
\code{topica.topic\_diversity} measure distinctness, \code{bootstrap\_stability}
checks that topics persist across seeds, and \code{word\_intrusion} and
\code{document\_intrusion} build the human-validation tasks of
\citet{chang2009reading}.

\subsection{Measure effects with honest uncertainty}

When the question is about variation across groups or over time, the result is a
regression of topic prevalence on covariates. \pkg{topica} implements the method
of composition \citep{roberts2016model}, which propagates topic-estimation
uncertainty into the effect, with cluster-robust standard errors and generalized
linear model links:

\begin{CodeChunk}
\begin{CodeInput}
from topica import stm

X, names = topica.one_hot(party)
model = topica.STM(num_topics=20, seed=42)
model.fit(docs, prevalence=X, prevalence_names=names)

draws   = stm.posterior_theta_samples(model, nsims=50, seed=0)
effects = stm.estimate_effect(draws, X, feature_names=names, cluster=source_id)
\end{CodeInput}
\end{CodeChunk}

The model-agnostic \code{topica.standard\_errors} extends this to any model and is
explicit about the kind of uncertainty it reports. For a clustering of embeddings,
which has no posterior over $\theta$, it declines to report confidence intervals
and says so, rather than returning a number that looks like one.

\subsection{Visualization and reporting}

The same contract drives a visualization layer, \code{topica.viz}, that turns a
fitted model into publication figures. A single call to \code{topica.plot\_report}
(or the interactive \code{dashboard}) assembles a composite report from
\code{topic\_word} and \code{doc\_topic} alone, so the same report renders for
every model in the family; Figure~\ref{fig:report} shows it for the poliblog
\code{STM}. The toolkit is built to be an \emph{honest} successor to \pkg{LDAvis}
\citep{sievert2014ldavis}: it draws confidence bands only where a posterior
exists, labels each topic-word weight for what it is, reports a closure-corrected
topic correlation rather than the spuriously strong raw version, flags \emph{dead}
and near-duplicate topics, and refuses panels a model cannot honestly support---a
clustering of embeddings with no posterior over $\theta$ gets no uncertainty
ribbons. Interactive views render through \pkg{Plotly}, and
\code{topica.prepare\_pyldavis} still feeds the original \pkg{LDAvis} when wanted.

\begin{figure}[t!]
\centering
\includegraphics[width=\textwidth]{fig_poliblog_report}
\caption{The one-call report \code{topica.plot\_report(model, texts, groups)} for
the poliblog structural topic model ($K = 15$): topics ordered by prevalence with
FREX labels, the coherence--exclusivity quality frontier (point area is
prevalence), the topic-correlation matrix, and topic prevalence by blog rating.
Every panel reads from the shared \code{topic\_word}\,/\,\code{doc\_topic} surface,
so the same report renders for any model in Table~\ref{tab:models}. Reproduced by
\code{paper/replication.py}.}
\label{fig:report}
\end{figure}

%% ===========================================================================
\section{Validation against reference implementations}\label{sec:validation}
%% ===========================================================================

A framework that reimplements published methods earns trust only by agreeing with
the implementations social scientists already use---so rather than assert that
agreement, we ship the scripts that run the reference packages themselves and
report what they produce. The checks in \code{parity/} drive Java \pkg{MALLET} and
the \proglang{R} \pkg{stm} and \pkg{keyATM} packages on the same tokenized inputs
as \pkg{topica}, align the fitted topics, and report the overlap;
\code{paper/replication.py} runs them end to end and skips, visibly, any whose
toolchain is absent. Because the engines share no code and no random number
generator, ``agreement'' means \emph{statistical} agreement on the fitted topics,
benchmarked against each reference's own run-to-run reproducibility---the only
honest yardstick for stochastic inference. We report it without inflating it.

\paragraph{Collapsed-Gibbs LDA.} \pkg{topica}'s LDA binds \pkg{RustMallet}, David
Mimno's \proglang{Rust} port of \pkg{MALLET}'s SparseLDA sampler
\citep{yao2009sparselda}, and reproduces its \code{train} command-line output
byte-for-byte: the topic-word and document-topic matrices are bit-identical to the
reference binary's (\code{tests/test\_cli\_parity.py}). Against the original
\proglang{Java} \pkg{MALLET}, whose independent random number generator rules out
byte equality, \pkg{topica} instead recovers the \emph{same} topics, aligning at a
top-word Jaccard and cosine of $1.000$ on a planted-topic corpus
(\code{parity/mallet\_parity.py}); \code{LabeledLDA} and \code{DMR} likewise match
\pkg{MALLET} at a topic cosine of $1.000$.

\paragraph{Keyword-assisted topics.} \pkg{topica}'s \code{KeyATM} shares the
\proglang{R} \pkg{keyATM} package's per-sweep asymmetric-$\alpha$ estimation and
\code{model\_fit} log-likelihood \citep{eshima2024keyatm}. \pkg{keyATM}
initializes at random, so even two \proglang{R} runs differ; on the poliblog
corpus the keyword topics---the ones the method pins---agree between \pkg{topica}
and \proglang{R} at a cosine of $0.91$, as close as \proglang{R}'s two seeds agree
with each other ($0.90$) (\code{parity/keyatm\_r\_compare.py}).

\paragraph{Structural topic model.} The structural topic model is non-convex, so
the right question is not whether two engines produce identical matrices---they
need not---but whether \pkg{topica} lands on \proglang{R}'s solution as closely as
\proglang{R} lands on its own. On the poliblog vignette, with both engines given
the same Spectral initialization, \pkg{topica}'s prevalence fit agrees with
\proglang{R} \pkg{stm}'s at a topic-word cosine of $0.973$ (per-topic median
$0.993$), far inside \proglang{R}'s own spectral-versus-random spread of about
$0.61$---\pkg{topica} reproduces \proglang{R}'s solution more closely than
\proglang{R}, reinitialized, reproduces itself (\code{parity/stm\_poliblog\_compare.py}).
For the content model the comparison is sharper still: the per-group topic-word
distributions agree at a cosine of $0.999$--$1.000$ (\code{parity/stm\_content\_r\_compare.py}).
The parameters replicate, and so does the conclusion: on the canonical vignettes
\pkg{topica} recovers the same prevalence effects the package documents, with
honest standard errors from the method of composition.

%% ===========================================================================
\section{Performance}\label{sec:performance}
%% ===========================================================================

We benchmark fit time only (construction and import excluded) on fixed-seed
synthetic corpora, on one 14-core Apple M-series machine, against the reference
implementations on matched iterations from a common initialization, so the
comparison measures per-iteration cost rather than time to convergence. Speed
depends on corpus size, vocabulary, the number of topics, and hardware, so the
numbers are orders of magnitude, not guarantees. Every comparison here is
reproduced by \code{benchmarks/speed\_vs\_r.py}, which times the identical fits in
\proglang{R} and \pkg{topica} (and is driven, with the validation checks, by
\code{paper/replication.py}).

\begin{table}[t!]
\centering
\caption{\pkg{topica} versus \proglang{R} \pkg{stm}, fit time on matched EM
iterations from a spectral start. \proglang{R} \pkg{stm} is single-threaded by
design; \pkg{topica} parallelizes the variational E-step.}
\label{tab:stm}
\small
\begin{tabular}{rrrrrrr}
\toprule
\multicolumn{3}{c}{Corpus} & \multicolumn{2}{c}{Single core} & \multicolumn{2}{c}{All cores}\\
\cmidrule(lr){1-3}\cmidrule(lr){4-5}\cmidrule(lr){6-7}
docs & vocab & $K$ & \pkg{topica} & \proglang{R} \pkg{stm} & \pkg{topica} & speedup \\
\midrule
1{,}000 & 500   & 10 & 0.50s & 3.03s & 0.14s & 22.5$\times$ \\
2{,}000 & 2{,}000 & 10 & 1.44s & 6.51s & 0.49s & 13.5$\times$ \\
5{,}000 & 5{,}000 & 20 & 8.97s & 26.3s & 2.75s & 9.8$\times$ \\
\bottomrule
\end{tabular}
\end{table}

The structural topic model is the comparison that matters most for social
scientists, since \proglang{R} \pkg{stm} is the field standard and its fits can
take minutes. Table~\ref{tab:stm} shows \pkg{topica} fitting the same model three
to six times faster on a single core and ten to twenty-two times faster across
cores. The keyword-assisted model matches \proglang{R} \pkg{keyATM}'s
\proglang{C++} sampler single-threaded (25.9s versus 24.5s on a
2{,}000-document, 10-topic fit of 1{,}000 sweeps) and runs about twice as fast on
four cores (12.1s), through a document-partitioned parallel sweep that the
\proglang{R} package has no equivalent of.

We report the LDA comparison straight, because it does not favor us. Against
\pkg{MALLET}, the original \proglang{Java} sampler, fit time is close: the two are
roughly at parity single-threaded, \pkg{topica}'s document-partitioned parallelism
pulls ahead on four cores (about $2\times$), and \pkg{topica} adds no JVM startup. Against
\pkg{tomotopy} \citep{tomotopy}, a
\proglang{C++} library in the same performance tier, plain LDA is a wash:
\pkg{tomotopy}'s tighter inner loop is about twenty percent ahead single-threaded,
and \pkg{topica}'s document-partitioned parallelism pulls even or slightly ahead
across cores. For plain LDA the two are interchangeable on speed; \pkg{topica}'s
advantage is the structural-topic-model, effect-estimation, and diagnostics stack
built around the sampler, not raw LDA throughput.

%% ===========================================================================
\section{Worked example: many models, one workflow}\label{sec:example}
%% ===========================================================================

\pkg{topica}'s value shows in two directions at once on a single corpus: across
models, where the shared contract makes a many-model comparison nearly free, and
within one model, where the structural topic model answers a substantive question
with honest uncertainty. Both run on the political-blog corpus of the \pkg{stm}
literature \citep{roberts2016model}---weblog posts from the 2008 \proglang{U.S.}
presidential campaign, rated by the blog's political leaning. The full, runnable
script is \code{paper/replication.py}.

\subsection{Comparing models behind one contract}

Because every model exposes the same surface, scoring a heterogeneous set of them
is one loop, not four pipelines. Here four models from two families---three
count-based and one embedding-based---are fit on the same corpus and scored by the
same diagnostics:

\begin{CodeChunk}
\begin{CodeInput}
import topica

# docs: preprocessed posts; conservative: (D, 1) indicator column.
def score(m):                # the SAME diagnostics for every model
    return (topica.coherence(m, docs).mean(),
            topica.exclusivity(m).mean(),
            topica.topic_diversity(m))

lda = topica.LDA(num_topics=15, seed=1); lda.fit(docs, iters=500)
ctm = topica.CTM(num_topics=15, seed=1); ctm.fit(docs, iters=25)
stm_model = topica.STM(num_topics=15, seed=1)
stm_model.fit(docs, conservative,
              prevalence_names=["conservative"], iters=25)

# An embedding model joins the same comparison: bring vectors,
# fit, and score it too.
bert = topica.BERTopic(reducer="pca", min_cluster_size=10, seed=42)
bert.fit_transform(docs, embeddings)      # any document vectors

for name, m in [("LDA", lda), ("CTM", ctm),
                ("STM", stm_model), ("BERTopic", bert)]:
    print(name, m.num_topics, *("%.3f" % v for v in score(m)))
\end{CodeInput}
\end{CodeChunk}

Table~\ref{tab:spanning} collects the result. The count-based models resolve
fifteen granular topics; the embedding clustering resolves three broad, highly
coherent macro-topics. Which granularity is right is the researcher's call, not
the software's---and the point is that the comparison needed to make that call
costs no translation between tools, because all four models are scored through the
identical interface.

\begin{table}[t!]
\centering
\caption{Four models from two families, fit on the same 2{,}000-post poliblog
corpus and scored by the same three diagnostics in one loop (the code above).
Coherence is windowed $c_v$ \citep{lau2014coherence}; exclusivity and diversity
measure topic distinctness. The comparison costs no extra plumbing because every
model presents the same surface. Numbers reproduce from \code{paper/replication.py}.}
\label{tab:spanning}
\small
\begin{tabular}{llrrrr}
\toprule
Model & Family & Topics & Coherence ($c_v$) & Exclusivity & Diversity \\
\midrule
\code{LDA}      & count-based     & 15 & 0.358 & 0.514 & 0.683 \\
\code{CTM}      & count-based     & 15 & 0.351 & 0.446 & 0.579 \\
\code{STM}      & count-based     & 15 & 0.350 & 0.435 & 0.565 \\
\code{BERTopic} & embedding-based &  3 & 0.459 & 0.526 & 0.693 \\
\bottomrule
\end{tabular}
\end{table}

A comparison behind one contract need not stop at diagnostics. Two count-based
models from different families can be aligned topic to topic and their covariate
effects set side by side, each carried with the uncertainty its own posterior
admits:

\begin{CodeChunk}
\begin{CodeInput}
from topica import stm

# lda (collapsed-Gibbs) and stm_model (variational) were fit above on the
# same corpus, so their topic-word columns share a vocabulary. Match them.
pairs = topica.align_topics(lda, stm_model)   # [(lda_topic, stm_topic, dist), ...]

# One estimator, two families: it draws the posterior each admits --- Dirichlet
# theta for the Gibbs LDA, logistic-normal theta for the STM --- so the
# conservative effect carries method-of-composition uncertainty in both.
eff_lda = stm.estimate_effect(lda, conservative, corpus=docs,
                              feature_names=["conservative"], nsims=100)
eff_stm = stm.estimate_effect(stm_model, conservative,
                              feature_names=["conservative"], nsims=100)

# Compare the conservative slope on matched topics, each with its own 95% CI.
for i, j, _ in pairs:
    a, b = eff_lda[i], eff_stm[j]   # coef[1] = conservative; coef[0] = intercept
    print("%+.3f [%+.3f, %+.3f]   %+.3f [%+.3f, %+.3f]" % (
        a.coef[1], a.ci_low[1], a.ci_high[1],
        b.coef[1], b.ci_low[1], b.ci_high[1]))
\end{CodeInput}
\end{CodeChunk}

The match is one assignment problem on the topic-word distances, and the single
\code{estimate\_effect} call serves both families, drawing Dirichlet $\theta$ for
the collapsed-Gibbs LDA and logistic-normal $\theta$ for the variational STM, so
each conservative effect arrives with a method-of-composition interval rather than
a bare point. Where two model families agree, within their intervals, on which
topics are partisan, that agreement is a cross-model robustness check no
single-model tool can pose; and the same call, asked of the embedding clustering
of Table~\ref{tab:spanning}, finds no posterior over $\theta$ and declines the
interval rather than manufacture one. This is the measurement stance of
Section~\ref{sec:design} made operational: uncertainty is propagated from topic
estimation into the effect where a posterior exists, and withheld where it does
not, uniformly across the model family.

\subsection{One model in depth: covariate effects with honest uncertainty}

The same corpus also carries a substantive question the structural topic model is
built for: how does topic prevalence differ by the blog's political rating? The
fit, FREX labeling, and effect plot are a handful of lines:

\begin{CodeChunk}
\begin{CodeInput}
from topica import Corpus, stm
import topica.viz as viz

corpus = Corpus.from_documents(
    docs, min_doc_freq=10, max_doc_fraction=0.5, rm_top=20)
model = topica.STM(num_topics=15, seed=1)
model.fit(docs, conservative,
          prevalence_names=["conservative"], iters=25)

labels = stm.label_topics(model.topic_word, model.vocabulary, n=3)
panel = viz.effect_plot(model, corpus, X=conservative,
                        feature_names=["conservative"], nsims=100)
panel.to_png("fig_poliblog_effect.pdf")
\end{CodeInput}
\end{CodeChunk}

\begin{figure}[t!]
\centering
\includegraphics[width=0.82\textwidth]{fig_poliblog_effect}
\caption{The results figure produced by the code above, on 2{,}000 posts and a
2{,}612-word vocabulary. Each row is one of $K=15$ topics, labeled by its top
three FREX words \citep{bischof2012frex}, ordered by the estimated difference in
prevalence between conservative and liberal blogs, with 95\% method-of-composition
intervals. Topics whose interval excludes zero are partisan; those that straddle
it are discussed at similar rates on both sides.}
\label{fig:poliblog}
\end{figure}

The fit recovers the partisan division the \pkg{stm} poliblog analysis documents
(Figure~\ref{fig:poliblog}). Conservative blogs disproportionately discuss the
Reverend Wright and Bill Ayers attacks on Obama
($+0.046$, 95\% CI $[0.034, 0.057]$), media coverage and bias ($+0.043$,
$[0.027, 0.059]$), and Israel ($+0.056$, $[0.044, 0.068]$); liberal blogs
disproportionately discuss McCain ($-0.084$, $[-0.094, -0.074]$) and the Bush
administration and torture ($-0.065$, $[-0.077, -0.053]$). Topics such as the
economy and abortion straddle zero, discussed at similar rates on both sides. The
intervals come from the method of composition, so they reflect the uncertainty in
estimating $\theta$, not only sampling variation in the regression. The full,
runnable version of this analysis, including the choice of $K$ and the
clustered-by-blog standard errors, is the script that produced this figure
(\code{paper/replication.py}) and the corresponding worked example in the
package documentation.

%% ===========================================================================
\section{Discussion}\label{sec:discussion}
%% ===========================================================================

\pkg{topica} unifies a fragmented toolchain, but it does not dissolve the
assumptions the underlying models carry. The count-based models are bag-of-words
models and inherit the limits of that representation. The embedding-based models
are only as good as the embeddings the user supplies, and the optional UMAP
reducer is not seedable in its current upstream form; \pkg{topica} contains that
nondeterminism by following the \pkg{BERTopic} pattern of a stochastic discovery
fit and a deterministic prediction phase, and warns when the reducer is used. The
conservative refusal to report intervals for models without a posterior is a
deliberate trade: it costs a number the analyst might have wanted and buys against
a number that would have been misleading.

The design commitment that the researcher owns the theoretical decisions is also
a limitation by choice. \pkg{topica} will not pick $K$ from a metric, will not
auto-label topics as if the labels were facts, and will not present a finished
result that looks authoritative and is theoretically ungrounded. A study done
well with \pkg{topica} leaves the researcher having made every decision a reviewer
will question, with the diagnostics to defend each one. As large-language-model
agents increasingly drive analysis pipelines, this division of labor matters more,
not less, and the package ships a companion guide that encodes it for agent use.

Future work includes amortized variational inference for the embedded topic model
to scale it past per-document EM, additional models, and a wider set of links and
estimands for the effect-estimation tools.

%% ===========================================================================
\section{Availability}\label{sec:availability}
%% ===========================================================================

\pkg{topica} is released under the Apache-2.0 license and installs from the
Python Package Index with \code{pip install topica}; the core requires only
\pkg{NumPy}. Source code, documentation, and worked examples are available at
\url{https://github.com/nealcaren/topica} and
\url{https://nealcaren.github.io/topica/}. A single script,
\code{paper/replication.py}, regenerates every quantitative artifact in this
article---the worked-example figure and table (Section~\ref{sec:example}) and,
where the reference toolchains are installed, the cross-implementation validation
(Section~\ref{sec:validation}) and speed comparisons (Section~\ref{sec:performance})
by running \pkg{MALLET} and the \proglang{R} \pkg{stm} and \pkg{keyATM} packages
themselves. \pkg{topica} reimplements published
methods and is validated against their reference implementations; it does not
replace the original work, and users should cite both \pkg{topica}
\citep{caren2026topica} and the model they use.

\bibliography{topica}

\end{document}
